\begin{figure}[tbp]
\centering
\begin{tikzpicture}[x=1cm,y=1cm,
 parent/.style={draw,rounded corners=2pt,minimum width=1.65cm,minimum height=.65cm,fill=black!4},
 child/.style={draw=blue!45!black,rounded corners=2pt,minimum width=3.25cm,minimum height=.65cm,fill=blue!4,font=\small},
 >={Stealth[length=4pt]}]
\node[anchor=east,font=\small] at (-.1,1.2) {Rows};
\node[parent] (rm) at (1,1.2) {$R_i^-(p)$};
\node[child] (r0) at (4.1,1.2) {subtree of $2p$};
\node[child] (r1) at (8.0,1.2) {subtree of $2p+1$};
\node[parent] (rp) at (11.1,1.2) {$R_i^+(p)$};
\draw[->] (rm.east)--(r0.west);
\draw[->] (r0.east)--(r1.west);
\draw[->] (r1.east)--(rp.west);
\node[anchor=east,font=\small] at (-.1,0) {Columns};
\node[parent] (cp) at (1,0) {$C_i^+(p)$};
\node[child] (c0) at (4.1,0) {subtree of $2p$};
\node[child] (c1) at (8.0,0) {subtree of $2p+1$};
\node[parent] (cm) at (11.1,0) {$C_i^-(p)$};
\draw[->] (cp.east)--(c0.west);
\draw[->] (c0.east)--(c1.west);
\draw[->] (c1.east)--(cm.west);
\end{tikzpicture}
\caption{The two orders at an internal node $(i,p)$, where $i<h$.
Each blue box contains the entire corresponding subtree on that axis.
At a leaf, the minus and plus names denote the same single position.
The reversed parent signs give the four comparisons in
\eqref{eq:tree-orders}.}
\label{fig:tree-orders}
\end{figure}
